% SOICT 2026 — Springer LNCS template (English)
% Based on: Latex-Template-for-Springer/samplepaper.tex
%
% Compile: ./compile.sh
%
\documentclass[runningheads]{llncs}

\usepackage{fontspec}
\usepackage{polyglossia}
\setdefaultlanguage{english}
\setotherlanguage{vietnamese}
% Shared XeLaTeX font setup for Vietnamese manuscripts.
% TeX Gyre Termes is a Times-compatible font that ships with TeX Live / MacTeX,
% covers Vietnamese, and needs no MS-Office fonts -> builds on any machine.
% If you have "Times New Roman" installed and your venue requires it exactly,
% swap the line below for: \setmainfont{Times New Roman}[Ligatures=TeX]
\setmainfont{TeX Gyre Termes}[Ligatures=TeX]

\usepackage{graphicx}
\usepackage{booktabs}
\usepackage{amsmath}

\graphicspath{{../../../results/ml_07_cross_method_comparison/}}

\begin{document}

\title{A Distributed Real-Time Intrusion Detection System on Jetson Nano Edge Cluster using Kafka and PySpark}

\titlerunning{Distributed Edge IDS on Jetson Nano}

\author{Nguyễn Vũ Thái\inst{1} \and
Bùi Văn Dũng\inst{2,3}}

\authorrunning{N. V. Thái et al.}

\institute{University of Transport and Communications, Ho Chi Minh City Campus, Vietnam\\
\email{nvthai@utc2.edu.vn}
\and
Vietnam National University Ho Chi Minh City, Vietnam
\and
University of Information Technology, Vietnam\\
\email{dungbv.17@grad.uit.edu.vn}}

\maketitle

\begin{abstract}
Deploying machine learning-based intrusion detection on resource-constrained edge devices requires balancing accuracy, latency, and operational cost.
We present a split-deployment IDS architecture across two NVIDIA Jetson Nano devices connected through Apache Kafka~\cite{kreps2011kafka}.
The system implements a two-stage pipeline: (1) a lightweight autoencoder anomaly gate on Jetson~\#1 filters benign traffic, and (2) a PySpark \texttt{PipelineModel} classifier on Jetson~\#2 processes only suspicious flows.
Central infrastructure (Kafka, PostgreSQL, InfluxDB, Grafana) runs on a host PC.
We evaluate three distributed modes---pipeline split, horizontal Kafka consumer scaling, and Spark standalone cluster---streaming CICIDS2017 network flows~\cite{sharafaldin2018toward}.
Experimental results show that the anomaly gate reduces Spark inference load while preserving detection accuracy, with sub-100\,ms batch latency on Jetson Nano under moderate load.

\keywords{Intrusion detection \and Edge computing \and Jetson Nano \and Kafka \and PySpark.}
\end{abstract}

\section{Introduction}

Network intrusion detection systems (IDS) must operate closer to data sources to reduce bandwidth and response time~\cite{shi2016edge}.
Edge devices such as the NVIDIA Jetson Nano offer limited CPU, memory, and thermal headroom, making monolithic ML inference challenging.

\textbf{Contributions:}
(1) A two-stage distributed IDS with anomaly gating and Spark-based classification;
(2) Three deployment modes for a two-node Jetson cluster;
(3) End-to-end performance evaluation including throughput, latency, and per-node resource usage;
(4) Open reproduction scripts integrated with our Spark ML training pipeline.

\section{Related Work}

Edge computing pushes analytics closer to data sources to reduce latency and bandwidth~\cite{shi2016edge,shi2021edge,khanna2023edge}.
Prior edge IDS studies~\cite{ten2019edge} benchmark single-node inference on resource-constrained boards but rarely evaluate multi-node Kafka pipelines with explicit role separation.
Streaming platforms such as Kafka~\cite{kreps2011kafka} enable decoupled producers and consumers; however, most ML-based IDS papers~\cite{khraisat2019survey,ferrag2020deep,vinayakumar2019deep} report offline batch accuracy without end-to-end throughput or per-node resource traces.

\textbf{Split inference.}
Anomaly detection surveys~\cite{chandola2009anomaly,ruff2021unifying} motivate two-stage designs: a lightweight gate followed by a heavier classifier.
Autoencoder-based detectors~\cite{lopez2020autoencoder} and federated edge IDS~\cite{mohammadi2023federated} explore similar ideas on servers or homogeneous edge nodes, but without PySpark MLlib classifiers on NVIDIA Jetson Nano~\cite{nvidia2019jetson}.
Our companion ML study selects SHAP Top-30 features on CICIDS2017~\cite{sharafaldin2018toward,lundberg2017unified}; this paper focuses on \emph{how} to deploy that model across two Jetsons with measurable gate skip ratio and three scaling modes.

Table~\ref{tab:related-work} positions our system against representative related work.

% English comparison table — \input from SOICT main.tex
\begin{table}[htbp]
  \caption{Comparison with related edge and streaming IDS deployments.}
  \label{tab:related-work}
  \centering
  \small
  \begin{tabular}{|p{2.2cm}|p{2.5cm}|p{1.5cm}|p{2.0cm}|p{2.5cm}|}
    \hline
    \textbf{Work} & \textbf{Focus} & \textbf{Platform} & \textbf{Streaming} & \textbf{Our difference} \\
    \hline
    Shi et al.~\cite{shi2016edge} & Edge vision & Generic & --- & Kafka + PySpark IDS on Jetson \\
    Ten et al.~\cite{ten2019edge} & Edge ML perf. & RPi / edge & Batch & End-to-end Kafka replay + 3 deploy modes \\
    Khanna et al.~\cite{khanna2023edge} & Edge IoT survey & Generic & Partial & Measured 2-node Jetson cluster \\
    Mohammadi et al.~\cite{mohammadi2023federated} & Federated IDS & Edge nodes & --- & Split anomaly gate + classifier \\
    Kreps et al.~\cite{kreps2011kafka} & Log streaming & Server & Kafka & Suspicious-flow topic forwarding \\
    Al-Rawashdeh \& Purdy~\cite{alrawashdeh2016toward} & Online RF IDS & Cloud VM & Near-line & Lightweight AE gate on Nano \\
    L{\'o}pez-Mart{\'i}n et al.~\cite{lopez2020autoencoder} & AE detection & GPU server & Offline & AE on gate, Spark only on suspicious \\
    NVIDIA Jetson~\cite{nvidia2019jetson} & Embedded GPU & Jetson Nano & --- & Role-based pipelines + Influx metrics \\
    \textbf{This work} & \textbf{Distributed edge IDS} & \textbf{2$\times$ Jetson} & \textbf{Kafka} & \textbf{Modes A/B/C + open scripts} \\
    \hline
  \end{tabular}
\end{table}


\section{System Architecture}

\subsection{Split Deployment Overview}

A host PC runs Kafka, PostgreSQL, InfluxDB, and Grafana. Two Jetson Nano nodes consume network-flow messages from topic \texttt{ids-network-flow}.
In pipeline-split mode, Jetson~\#1 forwards suspicious flows to \texttt{ids-suspicious-flow} for classification on Jetson~\#2.

\subsection{Pipeline Roles}

\begin{itemize}
  \item \textbf{anomaly\_gate:} sklearn autoencoder + scaler; MSE threshold from Exp~8.
  \item \textbf{classifier:} PySpark \texttt{PipelineModel} with 30 SHAP-selected features.
  \item \textbf{full:} combined pipeline for horizontal scaling via Kafka consumer groups.
\end{itemize}

\subsection{Monitoring and Alerting}

Each node tags metrics with \texttt{EDGE\_NODE\_ID} in InfluxDB and PostgreSQL. Alerts are sent from a designated primary node only.

\section{Implementation}

Training and model export run on a PC using Apache Spark ML~\cite{zaharia2010spark}.
The edge stack is implemented in Python (\texttt{raspberry/edge/}) with configurable roles via environment variables.

\section{Experimental Setup}

\subsection{Hardware and Software}

2$\times$ Jetson Nano 4GB, host PC (Docker Compose), WiFi LAN, PySpark 3.4.1, Kafka 7.5.

\subsection{Workloads}

CSV replay through \texttt{data\_sender.py} at 50--200 flows/s. Metrics: throughput, p50/p95 latency, CPU/RAM/temperature, attack detection rate, gate skip ratio.

\subsection{Baselines}

Single-node full pipeline vs.\ distributed modes A (split), B (horizontal), C (Spark cluster).

\section{Results}

Run \texttt{papers/soict2026/run\_benchmarks.sh} on Jetson hardware (modes \texttt{local}, \texttt{run}, \texttt{merge}).
End-to-end distributed metrics (throughput, p95 latency, gate skip \%) are exported to \texttt{results/benchmarks/summary.csv}.

\begin{table}[htbp]
  \caption{Distributed deployment comparison (fill from \texttt{summary.csv} after benchmarks).}
  \label{tab:benchmark}
  \centering
  \begin{tabular}{l c c c c}
    \toprule
    Mode & Throughput (rps) & Latency p95 (ms) & Gate skip (\%) & Attack F1 \\
    \midrule
    Single node & -- & -- & -- & -- \\
    A: Pipeline split & -- & -- & -- & -- \\
    B: Horizontal & -- & -- & N/A & -- \\
    C: Spark cluster & -- & -- & N/A & -- \\
    \bottomrule
  \end{tabular}
\end{table}

\begin{figure}[htbp]
  \centering
  \fbox{\parbox{0.9\textwidth}{\centering\vspace{2cm}[Insert architecture diagram: figures/architecture.pdf]\vspace{2cm}}}
  \caption{Split-deployment architecture (PC + 2$\times$ Jetson Nano).}
  \label{fig:architecture}
\end{figure}

\section{Discussion}

Pipeline split offloads Spark from the majority of benign traffic. Horizontal scaling improves throughput when Kafka partitions $\geq$ 2. Spark cluster mode adds coordination overhead on Nano-class hardware.

Model selection (Decision Tree / SHAP Top-30) follows our companion ML study on CICIDS2017.

\section{Conclusion}

We presented a reproducible distributed edge IDS for Jetson Nano clusters with Kafka streaming and PySpark inference.
Future work includes TensorRT optimization and federated retraining under drift.

\begin{credits}
\subsubsection{\ackname}
This work was conducted at the University of Transport and Communications (Ho Chi Minh City Campus)
and the University of Information Technology, Vietnam National University Ho Chi Minh City,
as part of a graduation thesis research project.

\subsubsection{\discintname}
The authors have no competing interests to declare relevant to this article.
\end{credits}

\bibliographystyle{../../latex/splncs04}
\bibliography{references,../../latex/related_work}

\end{document}
